\section{Связанные работы}
\label{sec:related}

Численные различия между реализациями одной операции рассматриваются
в анализе ошибок вычислений с плавающей точкой. Обратный анализ
ошибок~\cite{Wilkinson1963,Higham2002} и описание округления~\cite{Goldberg1991}
объясняют, почему изменение порядка суммирования может менять результат
при неизменных входах. Для оценки накопления таких отклонений используются
и вероятностные модели~\cite{HighamMary2019}. Их применение требует
предположений о средних и взаимной зависимости ошибок, которые нельзя
установить по одному воспроизводимому запуску.

В симуляторе к арифметике отдельной операции добавляется вопрос
о том, какие именно операции исполняются. Модели матричных блоков
MMA-Sim и работа Khattak и Mikaitis~\cite{MMASim2025,Khattak2025}
рассматривают особенности аппаратной арифметики; SCALE-Sim
и ASTRA-Sim~\cite{Samajdar2020,Rashidi2020} предназначены для
архитектурного моделирования. Компиляторы OpenXLA, TVM и
MLIR~\cite{OpenXLA2023,Chen2018TVM,Lattner2021MLIR}, в свою очередь,
преобразуют описание сети до исполнения вычислительными блоками.
Поэтому при исследовании численных результатов нужно различать замену
реализации матричной операции и изменение всего пути от описания модели
до её исполнения.

После выделения локальных расхождений возникает вопрос об их переносе
через остальные слои. Для трансформеров он связан, в частности, с
нормализацией и остаточными связями, складывающими вход блока
с результатом его преобразования. LayerNorm и
RMSNorm~\cite{Ba2016,Zhang2019} используют статистики текущего входа,
а расположение нормализации влияет на градиенты~\cite{Xiong2020,Wang2022}.
Это обосновывает сравнение разных схем блоков, но само по себе
не устанавливает сжатия любого возмущения. Аналогично исследования
обучения со смешанной разрядностью~\cite{Micikevicius2018,Wang2018FP8}
связывают арифметику с качеством в конкретных режимах; они не задают
универсального порога относительной ошибки тензора для любой модели.

В настоящей работе эти уровни рассматриваются последовательно:
сначала сравниваются пути исполнения и изолированные слои, затем
измеряется перенос отклонений к выходам и градиентам компонента сети.
Такой порядок позволяет сопоставить локальные и итоговые расхождения,
не предполагая заранее одинакового статистического механизма всех ошибок.
